\section{Multi-dataset Validation}
\label{sec:multidataset}

The original two-dataset evaluation (SQuAD + PubMedQA) was deliberately
designed to contrast a domain-matched setting against a
domain-mismatched one, but is inadequate for testing whether the
coherence paradox generalizes across question styles, document genres,
and reasoning structures. This section reports an extension to six
datasets covering open-domain factoid QA, multi-hop reasoning, and a
specialized financial domain.

\subsection{Datasets}

\begin{description}
  \item[\textsc{SQuAD}~\citep{rajpurkar2016squad}] reading comprehension
        on Wikipedia paragraphs. Domain-matched single-document QA.
  \item[\textsc{PubMedQA}~\citep{jin2019pubmedqa}] biomedical research
        QA over abstracts. Domain-mismatched single-document QA.
  \item[\textsc{NaturalQuestions}~\citep{kwiatkowski2019nq}] real Google
        queries paired with Wikipedia long-answer paragraphs. Tests
        open-domain factoid retrieval.
  \item[\textsc{TriviaQA}~\citep{joshi2017triviaqa}] trivia questions
        paired with Wikipedia/web evidence. Tests retrieval over
        heterogeneous evidence formats.
  \item[\textsc{HotpotQA}~\citep{yang2018hotpotqa}] multi-hop questions
        requiring reasoning over $\geq$2 paragraphs (distractor split).
        Tests whether the paradox manifests under multi-hop pressure.
  \item[\textsc{FinanceBench}~\citep{islam2023financebench}] open-book
        questions over corporate 10-K filings. Tests a specialized,
        long-document, professional domain.
\end{description}

The complete loader inventory is in \texttt{src/dataset\_loaders.py}.
All loaders share a common \texttt{(documents, qa\_pairs)} interface so
the multi-dataset harness can iterate without per-dataset branching, and
all use a fixed seed (42) for sample selection.

\subsection{Coherence-paradox table}

\begin{table}[!htbp]
\centering
\caption{Coherence paradox at six-dataset scale. Filled from
\texttt{results/multidataset/summary.csv}. \texttt{paradox\_drop} =
faithfulness(baseline) $-$ faithfulness(hcpc\_v1); \texttt{v2\_recovery}
= faithfulness(hcpc\_v2) $-$ faithfulness(hcpc\_v1). Positive paradox\_drop
confirms the failure mode; positive v2\_recovery confirms HCPC-v2's
selective gating helps.}
\label{tab:multidataset}
\small
\begin{tabular}{lcccccc}
\toprule
Dataset & $n$ & Faith\textsubscript{base} & Faith\textsubscript{v1} &
        Faith\textsubscript{v2} & paradox\_drop & v2\_recovery \\
\midrule
SQuAD          & TBD & TBD & TBD & TBD & TBD & TBD \\
PubMedQA       & TBD & TBD & TBD & TBD & TBD & TBD \\
NaturalQs      & TBD & TBD & TBD & TBD & TBD & TBD \\
TriviaQA       & TBD & TBD & TBD & TBD & TBD & TBD \\
HotpotQA       & TBD & TBD & TBD & TBD & TBD & TBD \\
FinanceBench   & TBD & TBD & TBD & TBD & TBD & TBD \\
\bottomrule
\end{tabular}
\end{table}

\subsection{Generator stratification (Item~A3)}

The same multi-dataset query set is rerun across three generators
(\texttt{src/generators.py} registry):

\begin{itemize}
  \item \textbf{Mistral-7B-Instruct-v0.2} (Mistral AI). Primary generator.
  \item \textbf{Llama-3-8B-Instruct} (Meta). Cross-architecture replication.
  \item \textbf{Qwen2.5-7B-Instruct} (Alibaba). Tests whether the paradox
        is specific to Western-pretrained model families.
\end{itemize}

\begin{table}[!htbp]
\centering
\caption{Generator-stratified faithfulness on the multidataset query set
(values from \texttt{results/multidataset/by\_generator.csv}). The
paradox is \emph{generator-invariant} if the sign of paradox\_drop is
preserved across all three rows for each dataset.}
\label{tab:by_generator}
\small
\begin{tabular}{llccc}
\toprule
Dataset & Condition & Mistral & Llama-3 & Qwen2.5 \\
\midrule
SQuAD     & baseline / v1 / v2 & TBD & TBD & TBD \\
PubMedQA  & baseline / v1 / v2 & TBD & TBD & TBD \\
NaturalQs & baseline / v1 / v2 & TBD & TBD & TBD \\
TriviaQA  & baseline / v1 / v2 & TBD & TBD & TBD \\
HotpotQA  & baseline / v1 / v2 & TBD & TBD & TBD \\
FinanceB. & baseline / v1 / v2 & TBD & TBD & TBD \\
\bottomrule
\end{tabular}
\end{table}

\subsection{Reproducibility checkpoints}

The multi-dataset harness (\texttt{experiments/run\_multidataset\_validation.py})
is checkpointed at the (dataset, generator, condition) tuple level: each
completed tuple is appended to \texttt{completed\_tuples.json} with its
per-query CSV, and the runner re-enters by skipping completed tuples.
This makes it safe to re-run after partial failure (e.g., Self-RAG
unavailable on M4) without recomputing finished tuples.
